% Chapter 2 — Literature Review.
\section{Literature Review}

Hallucination detection is an active research area, and the existing methods fall into three broad families. The first family asks a language model to evaluate output consistency, the second trains neural classifiers over frozen encoders, and the third builds on hand-crafted features and interpretable models. This section reviews each family and then positions the present work within it.

LLM-based detection relies on model-internal signals or on a separate judge model. SelfCheckGPT \cite{selfcheck} samples several responses and treats their agreement as a hallucination signal, and HaluAgent \cite{haluagent} shows that small language models can act as cheaper detectors. Luna \cite{luna} takes the judge idea further with a lightweight evaluation model that matches larger judges at a fraction of the cost. The common limitation is that judge verdicts give no feature-level reason, which makes a single failure hard to debug.

A second line of work trains classifiers on top of pretrained encoders. A hybrid framework in \cite{ieeeTai} combines encoder features with lightweight models and reports strong results on hallucination benchmarks. These methods are accurate, yet they need the encoder weights, they run on GPUs, and their embeddings are not directly interpretable.

The third family keeps the pipeline interpretable by construction. A multi-signal framework in \cite{ijert} combines retrieval grounding, NLI verification, and an XGBoost classifier, and \cite{multimedia} scores hallucination with token-level SHAP-LIME attributions. The present work differs from \cite{multimedia} in an important way. It explains predictions at the level of semantically meaningful evidence-consistency features rather than raw tokens, and it adds calibrated risk scores and a cross-domain evaluation that both papers lack.

On the benchmark side, HaluEval \cite{haluEval} and TruthfulQA \cite{truthqa} are standard test beds, and FActScore \cite{factscore} measures factual precision in long-form generation. RAGTruth \cite{ragtruth} contributes naturally generated responses with human word-level annotations, which is closer to production data. FaithBench \cite{faithbench} shows that current detectors still struggle on challenging modern-LLM outputs, and SpikeScore \cite{spikescore} demonstrates that cross-domain generalization remains an open problem.

The positioning of this paper follows from these observations. Prior work has explored hallucination detection, feature engineering, calibration, and SHAP in isolation, and this work combines them into a compact, course-feasible, calibrated and explainable risk-prediction pipeline with a deployable interface.
